% Part: proof-theory
% Chapter: proof-search
% Section: introduction

\documentclass[../../../include/open-logic-section]{subfiles}

\begin{document}

\olfileid{pt}{ps}{int}

\olsection{Introduction}

The axioms and inference rules of !!{proof} systems like \Log{G3c} or
\Log{N1i} define what trees of sequents or !!{formula}s count as
correct !!{proof}s. Given such a tree of sequents or !!{formula}s
(possibly with additional information such as which rules are applied
at each step, or labels for assumptions and inference rules that
discharge them) the definitions of axioms and rules allow us to check
whether the given tree is a correct !!{proof}. It is a diferent, and
typically more difficult, problem to \emph{find} !!a{proof} of a given
sequent or of a given !!{formula} from given assumptions. That is the
problem of \emph{proof search}.

There is a very simple-minded procedure for proof search. We can think
of !!{formula}s as sequences of symbols from a finite alphabet.
Sequents and trees of sequents or !!{formula}s can be thought of as
sequences of !!{formula}s as well (perhaps using special brackets to
indicate tree branching). The axioms and rules allow us to
mechanically check if a sequence of symbols represents a correct
!!{proof}. Consequently, we can mechanically generate \emph{all
possible} correct !!{proof}s (by generating all sequences of symbols
from the alphabet used to represent !!{proof}s and checking for each
candidate sequence whether it represents a correct !!{proof}). The
simple-minded proof search precedure is: generate all correct
!!{proof}s until we find one that is a !!{proof} of the given sequent
(or of the given !!{formula} with open assumptions among the given
ones).

A better method is to use the naive method you'd use to find
!!a{proof} by hand: you start with the end-sequent, pick !!a{formula},
and apply an inference rule ``backwards'' to generate a premise or
premises which, if they had !!{proof}s, would combine with the last
inference to !!a{proof} of the given sequent. A similar method
(although more complicated) is what you'd use to find !!a{proof} in a
natural deduction system.

But this method is not fool-proof: if you pick the wrong rule, or
!!{formula}s in the wrong order, you may hit a dead end. It is also
not completely specified: at each point, there may be more than one
possible !!{formula} to chose, or more than one rule to apply
backwards. A \emph{proof search algorithm} is a completely specified
procedure that will find !!a{proof} if one exists (i.e., is
fool-proof).

Some !!{proof} systems lend themselves to simpler proof search
algorithms than others. In sequent systems, the !!{main operator} of
!!a{formula}, and whether it occurs on the left or the right,
determines which rule to apply backwards. E.g., if you want to find
!!a{proof} of $\Gamma \Sequent \Delta, !A \land !B$ in which $!A \land
!B$ occurs on the right, you should apply \RightR{\land} backwards and
generate two corresponding premises $\Gamma \Sequent \Delta, !A$ and
$\Gamma \Sequent \Delta, !B$. If your system is \Log{G1c}, however,
the presence of contraction makes things more difficult because it
gives you more options: You could also apply \RightR{\Contraction}
backward and generate the premise $\Gamma \Sequent \Delta, !A \land
!B, !A \land !B$ instead. And then the premise $\Gamma \Sequent
\Delta, !A \land !B, !A \land !B, !A \land !B$, and so on. In
\Log{G2c} it's even worse. Applying the rule \RightR{\land} backwards
also means you have to guess $\Gamma_1$, $\Gamma_2$ such that $\Gamma
= \Gamma_1 \cup \Gamma_2$ and $\Delta_1$, $\Delta_2$ such that $\Delta
= \Delta_1 \cup \Delta_2$ and try your luck with the premises
$\Gamma_1 \Sequent \Delta_1, !A$ and $\Gamma_2 \Sequent \Delta_2, !B$.
A wrong guess might lead you into a dead end later on, and you'd have
to backtrack to the beginning and pick different $\Gamma_1$,
$\Gamma_2$, $\Delta_1$, $\Delta_2$. If the !!{formula} is $!A \lor !B$
you have to chose one of the two possibilities for \RightR{\lor},
producing one of the two possible premises $\Gamma \Sequent \Delta,
!A$ or $\Gamma \Sequent \Delta, !B$. A wrong choice would again lead
to having to backtrack.

The system \Log{G3c} does not have these problems. It has no
structural rules, and the choice of inference rule is determined by
the !!{main operator} and the side the !!{formula} occurs on. Hence,
\Log{G3c} is better suited to proof search than \Log{G1c} or
\Log{G2c}. A successful search will yield !!a{proof} in \Log{G3c}, and
that !!{proof} can then be translated into \Log{G1c} or \Log{G2c} (or
\Log{N1c} for that matter). So having a proof search algorithm for
\Log{G3c} and translation algorithms for \Log{G3c} !!{proof}s to
!!{proof}s in other systems also solves the proof search methods for
those systems.

How do we know that a proof search algorithm is fool-proof? This would
involve showing that if the algorithm fails to find !!a{proof}, no
!!{proof} can exist. After all, if we pick a bad algorithm, we might
run into cases where it fails to find a proof even if one exists. This
is not a problem the simple minded algorithm has, because it searches
through all possible !!{proof}s. But a proof search algorithm is
supposed to be efficient and expend the minimum amount of effort.

The key idea to showing that a proof search algorithm is fool-proof is
to show that if the algorithm fails, we can show from the way it fails
that the sequent can't possibly have !!a{proof}. And the way to show
that is to show that it is not valid. In classical first-order logic,
valid sequents $\Gamma \Sequent \Delta$ are those where every
!!{structure} makes at least one !!{formula} in $\Gamma$ is false or
at least one !!{formula} in~$\Delta$ is true. \Log{G3c} is
\emph{sound}, i.e., it !!{prove}s only valid sequents. This is
established by induction on !!{proof}s: Axioms are obvioulsy valid,
and if the premises of a rule are valid, so is the conclusion. To show
that a proof search algorithm is fool-proof we show that if the search
for !!a{proof} of $\Gamma \Sequent \Delta$ fails, we can define
!!a{structure} in which all !!{formula}s in $\Gamma$ are true and all
!!{formula}s in~$\Delta$ are false. This also establishes that
\Log{G3c} is \emph{complete}, i.e., that if $\Gamma \Sequent \Delta$
is valid, there is a \Log{G3c}-!!{proof} of it. Since all failed
proof searches show that $\Gamma \Sequent \Delta$ is invalid, every
proof search for a valid sequent must succeed.

\end{document}
